\cotefile{134-4}
\scanpage{84}
\mspage{P 63}
\noindent ou la géométrie algébrique. (Ainsi, la notion de ``topos'' refuse d'entrer dans le ``sac bourbachique'' des structures, décidément étroit aux entournures~!)\ \ En se décidant, en pleine connaissance de cause, certes, à ne pas s'engager dans cette ``galère'', Bourbaki a par là-même renoncé à son ambition initiale, qui était de fournir\ \ \mbox{l e s}\ \ fondements et\ \ \mbox{l e}\ \ langage de base pour l'ensemble de la mathématique contemporaine.

\noindent\hspace*{2em}Il a,\ \ par contre, fixé un langage et, en même temps, un certain\ \ \mbox{s t y l e} d'écriture et d'approche de la mathématique. Ce style\ \ était à l'origine le reflet (très partiel) d'un certain\ \ \mbox{e s p r i t ,} vivant et direct héritage de Hilbert. Au cours des années cinquante et soixante, ce style a fini par s'imposer - pour le meilleur et (surtout) pour le pire. Depuis une vingtaine d'années, il a fini par devenir un rigide\ \ \mbox{`` c a n o n ''}\ \ d'une ``rigueur'' de\ \ \ \ \ \ pure façade, dont l'esprit qui l'animait jadis semble disparu sans retour.

\noindent\rule{0.15\linewidth}{0.45pt}\par\medskip

\noindent\hspace*{3em}Il semblerait bien qu'en tant que serviteur d'une vaste vision unificatrice née en moi, je sois ``unique en mon genre'' dans l'histoire de la mathématique de l'origine à nos jours. Désolé d'avoir l'air de vouloir me singulariser plus qu'il ne paraît permis~! A mon propre soulagement, je crois pourtant discerner une sorte de\ \ \mbox{f r è r e}\ \ potentiel (et providentiel~!). J'ai déjà eu tantôt l'occasion de l'évoquer, comme le premier dans la lignée de mes ``frères de tempérament''~: c'est\ \ \mbox{E v a r i s t e\ \ \ G a l o i s .}\ \ Dans sa courte et fulgurante vie (*), je crois discerner l'amorce d'une grande vision - celle justement des \handwritten{``}épousailles du nombre et de la grandeur'', dans une vision géométrique nouvelle. J'évoque ailleurs dans Récoltes et Semailles (**) comment, il y a deux ans, est apparu en moi cette intuition soudaine~: que dans le travail mathématique qui à ce moment exerçait sur moi la fascination la plus puissante, j'étais en train de ``reprendre l'héritage de Galois''. Cette intuition,\ \ rarement évoquée depuis, a pourtant eu le temps de mûrir en silence. La réflexion rétrospective sur mon oeuvre que je poursuis depuis trois semaines y aura sûrement encore contribué. La filiation la plus directe que je crois reconnaître à présent avec un mathématicien du passé, est bien celle qui me relie à Evariste Galois. A tort ou à raison, il me semble que cette vision que j'ai développée pendant quinze années de ma vie, et qui a continué encore à mûrir en moi et à s'enrichir pendant les seize années écoulées depuis mon départ de la scène mathématique - que cette vision est aussi celle que Galois n'aurait pu s'empêcher de développer (***), s'il s'était trouvé dans

\medskip\noindent\rule{0.15\linewidth}{0.45pt}

\noindent (*) Evariste Galois (1811-1832) est mort dans un duel, à l'âge de vingt\deleted{\unsure{-et-un}} ans. Il y a, je crois, plusieurs biographies de lui. J'ai lu comme jeune homme une biographie romancée, écrite par le physicien Infeld, qui m'avait beaucoup frappée à l'époque.

\noindent (**) Voir ``L'héritage de Galois'' (ReS I, section 7).

\noindent (***) Je suis persuadé d'ailleurs qu'un Galois serait allé bien plus loin encore que\ \ je n'ai été. D'une part à cause de ses dons tout à fait exceptionnels (que je n'ai pas reçus\ \ en partage, quant à moi). D'autre part parce qu'il est probable qu'il n'aurait pas, comme moi, laissé se distraire la majeure part de son énergie

% archive sheet number: 83
